\begin{bt}%[9D6H18-2]
% nguon: Tuyển tập đề tham khảo môn Toán Tuyển sinh 10 năm học 2025-2026 TP.HCM – Đề tham khảo Huyện Nhà Bè (Đề 1) – Bài 1
Cho hàm số $y=2x^2$ có đồ thị là Parabol $(P)$.
\begin{enumerate}[a)]
	\item Vẽ đồ thị $(P)$ trên hệ trục toạ độ $Oxy$.
	\item Xác định toạ độ của điểm $A$ biết điểm $A$ nằm trên đồ thị $(P)$ có hoành độ bằng $-\dfrac{1}{2}$.
\end{enumerate}
\loigiai{
\begin{enumerate}[a)]
	\item Bảng giá trị tương ứng của $x$ và $y$:
	\begin{center}
	\begin{tabular}{|c|c|c|c|c|c|}
	\hline
	$x$ & $-2$ & $-1$ & $0$ & $1$ & $2$ \\ \hline
	$y=2x^2$ & $8$ & $2$ & $0$ & $2$ & $8$ \\ \hline
	\end{tabular}
	\end{center}
	Đồ thị $(P)$ là một đường parabol đi qua gốc toạ độ $O(0;0)$, nhận trục tung $Oy$ làm trục đối xứng và quay bề lõm lên phía trên.
	\item Điểm $A$ thuộc $(P)$ có hoành độ $x_A = -\dfrac{1}{2}$.\\
	Thay $x_A = -\dfrac{1}{2}$ vào phương trình của $(P)$, ta được tung độ của điểm $A$:
	$$ y_A = 2 \cdot \left(-\dfrac{1}{2}\right)^2 = 2 \cdot \dfrac{1}{4} = \dfrac{1}{2}. $$
	Vậy toạ độ của điểm $A$ là $A\left(-\dfrac{1}{2}; \dfrac{1}{2}\right)$.
\end{enumerate}
}
\end{bt}

\begin{bt}%[9D6H20-2]
% nguon: Tuyển tập đề tham khảo môn Toán Tuyển sinh 10 năm học 2025-2026 TP.HCM – Đề tham khảo Huyện Nhà Bè (Đề 1) – Bài 2
Cho phương trình bậc hai: $x^2-4x+3=0$.
\begin{enumerate}[a)]
	\item Chứng minh phương trình trên luôn có hai nghiệm phân biệt $x_1, x_2$.
	\item Không giải phương trình, hãy tính giá trị của biểu thức:
	$$ A = \dfrac{5x_1-x_2}{x_1} - \dfrac{x_1-5x_2}{x_2}. $$
\end{enumerate}
\loigiai{
\begin{enumerate}[a)]
	\item Phương trình $x^2-4x+3=0$ có các hệ số $a=1$, $b'=-2$, $c=3$.\\
	Biệt thức thu gọn:
	$$ \Delta' = (-2)^2 - 1 \cdot 3 = 4 - 3 = 1 > 0. $$
	Vì $\Delta' > 0$ nên phương trình luôn có hai nghiệm phân biệt $x_1, x_2$.
	\item Theo hệ thức Vi-ét, ta có:
	$$ \begin{cases} S = x_1+x_2 = 4 \\ P = x_1 x_2 = 3. \end{cases} $$
	Quy đồng mẫu thức biểu thức $A$:
	$$ A = \dfrac{x_2(5x_1-x_2) - x_1(x_1-5x_2)}{x_1 x_2} = \dfrac{5x_1 x_2 - x_2^2 - x_1^2 + 5x_1 x_2}{x_1 x_2} = \dfrac{10x_1 x_2 - (x_1^2+x_2^2)}{x_1 x_2}. $$
	Biến đổi tử thức:
	$$ 10x_1 x_2 - (x_1^2+x_2^2) = 10x_1 x_2 - \left[(x_1+x_2)^2 - 2x_1 x_2\right] = 12x_1 x_2 - (x_1+x_2)^2 = 12P - S^2. $$
	Thay $S=4$ và $P=3$ vào biểu thức $A$, ta được:
	$$ A = \dfrac{12 \cdot 3 - 4^2}{3} = \dfrac{36 - 16}{3} = \dfrac{20}{3}. $$
\end{enumerate}
}
\end{bt}

\begin{bt}%[9D8H26-3]
% nguon: Tuyển tập đề tham khảo môn Toán Tuyển sinh 10 năm học 2025-2026 TP.HCM – Đề tham khảo Huyện Nhà Bè (Đề 1) – Bài 3
Biểu đồ cột kép biểu diễn số học sinh nam và số học sinh nữ của lớp 9A có sở thích chơi một số môn thể thao: Bóng đá (Nam: $12$, Nữ: $2$), Bóng rổ (Nam: $8$, Nữ: $6$), Bơi lội (Nam: $5$, Nữ: $10$). Biết rằng mỗi học sinh chỉ nêu một môn thể thao yêu thích nhất.
\begin{enumerate}[a)]
	\item Trong các môn thể thao trên, học sinh nam của lớp thích môn thể thao nào nhất?
	\item Chọn ngẫu nhiên một học sinh trong danh sách trên của lớp 9A, tính xác suất của các biến cố sau:
	\begin{itemize}
		\item Biến cố $A$: ``Học sinh được chọn là học sinh nam''.
		\item Biến cố $B$: ``Học sinh được chọn thích môn bóng rổ''.
		\item Biến cố $C$: ``Học sinh được chọn là nữ và không thích môn bơi lội''.
	\end{itemize}
\end{enumerate}
\loigiai{
\begin{enumerate}[a)]
	\item Dựa vào số liệu khảo sát, số lượng học sinh nam yêu thích các môn thể thao lần lượt là: Bóng đá ($12$ học sinh), Bóng rổ ($8$ học sinh), Bơi lội ($5$ học sinh).\\
	Vậy trong các môn thể thao trên, học sinh nam thích môn Bóng đá nhất.
	\item Tổng số học sinh của lớp 9A là:
	$$ n(\Omega) = (12 + 2) + (8 + 6) + (5 + 10) = 14 + 14 + 15 = 43\text{ (học sinh)}. $$
	\begin{itemize}
		\item Tổng số học sinh nam của lớp là $12 + 8 + 5 = 25$ học sinh.\\
		Xác suất của biến cố $A$ là:
		$$ P(A) = \dfrac{25}{43}. $$
		\item Tổng số học sinh thích môn bóng rổ (cả nam và nữ) là $8 + 6 = 14$ học sinh.\\
		Xác suất của biến cố $B$ là:
		$$ P(B) = \dfrac{14}{43}. $$
		\item Số học sinh nữ không thích môn bơi lội là số bạn nữ thích bóng đá hoặc bóng rổ:
		$$ 2 + 6 = 8\text{ (học sinh)}. $$
		Xác suất của biến cố $C$ là:
		$$ P(C) = \dfrac{8}{43}. $$
	\end{itemize}
\end{enumerate}
}
\end{bt}

\begin{bt}%[9D1H3-4]
% nguon: Tuyển tập đề tham khảo môn Toán Tuyển sinh 10 năm học 2025-2026 TP.HCM – Đề tham khảo Huyện Nhà Bè (Đề 1) – Bài 4
Bảng cước phí dịch vụ Mobicard (đã bao gồm thuế giá trị gia tăng VAT) quy định như sau: nếu thời lượng cuộc gọi trong $6$ giây đầu thì cước phí tính trọn gói là $118$ đồng; kể từ sau giây thứ $6$ trở đi, nhà mạng tính thêm $19{,}5$ đồng cho mỗi giây phát sinh.
\begin{enumerate}[a)]
	\item Hãy thiết lập công thức biểu diễn số tiền cước $y$ (đồng) phải trả khi gọi trên $6$ giây, với $x$ là số giây gọi tính từ sau giây thứ $6$ trở đi ($x \in \mathbb{N}^*$).
	\item Bạn Khang thực hiện một cuộc gọi và phải trả số tiền cước là $2419$ đồng. Hỏi bạn Khang đã gọi cuộc gọi đó trong thời gian bao lâu?
\end{enumerate}
\loigiai{
\begin{enumerate}[a)]
	\item Với $x$ là số giây gọi phát sinh kể từ sau giây thứ $6$, số tiền cước tính thêm là $19{,}5x$ (đồng).\\
	Cộng với mức cước ban đầu của $6$ giây đầu là $118$ đồng, công thức biểu diễn số tiền cước $y$ phải trả là:
	$$ y = 118 + 19{,}5x\text{ (đồng)}. $$
	\item Theo đề bài, tổng số tiền cước bạn Khang phải trả là $2419$ đồng, ta có phương trình:
	$$ 118 + 19{,}5x = 2419 \iff 19{,}5x = 2419 - 118 \iff 19{,}5x = 2301 \iff x = \dfrac{2301}{19{,}5} = 118\text{ (giây)}. $$
	Tổng thời gian của cuộc gọi bạn Khang đã thực hiện là:
	$$ 6 + 118 = 124\text{ (giây)}, $$
	tương đương $2$ phút $4$ giây.
\end{enumerate}
}
\end{bt}

\begin{bt}%[9H0V31-3]
% nguon: Tuyển tập đề tham khảo môn Toán Tuyển sinh 10 năm học 2025-2026 TP.HCM – Đề tham khảo Huyện Nhà Bè (Đề 1) – Bài 5
Một chiếc bình có phần thân dưới dạng hình trụ với đường kính đáy $8\text{ cm}$. Khi đặt bình thẳng đứng trên mặt phẳng nằm ngang, mực nước trong bình cao $10\text{ cm}$.
\begin{enumerate}[a)]
	\item Tính thể tích nước chứa trong bình (lấy $\pi \approx 3{,}14$, kết quả làm tròn đến hàng phần trăm của $\text{cm}^3$).
	\item Nếu dùng chiếc bình này khi chứa đầy nước rồi rót đầy vào các chiếc ly có dạng hình lập phương cạnh $4\text{ cm}$ (biết thể tích toàn phần của bình là $753{,}6\text{ cm}^3$) thì rót được tối đa bao nhiêu ly nước đầy như vậy?
\end{enumerate}
\loigiai{
\begin{enumerate}[a)]
	\item Bán kính đáy của phần hình trụ là $r = \dfrac{8}{2} = 4\text{ cm}$.\\
	Chiều cao của cột nước khi đặt bình thẳng đứng là $h = 10\text{ cm}$.\\
	Thể tích nước chứa trong bình là:
	$$ V_{\text{nước}} = \pi r^2 h = 3{,}14 \cdot 4^2 \cdot 10 = 502{,}4\text{ cm}^3. $$
	\item Thể tích của một chiếc ly hình lập phương cạnh $a = 4\text{ cm}$ là:
	$$ V_{\text{ly}} = a^3 = 4^3 = 64\text{ cm}^3. $$
	Khi bình chứa đầy nước với thể tích toàn phần là $753{,}6\text{ cm}^3$, số chiếc ly rót đầy tối đa được là:
	$$ N = \dfrac{753{,}6}{64} = 11{,}775. $$
	Vì số ly nước đầy phải là một số nguyên dương nên số ly rót đầy tối đa được là $11$ ly.
\end{enumerate}
}
\end{bt}

\begin{bt}%[9D1H3-4]
% nguon: Tuyển tập đề tham khảo môn Toán Tuyển sinh 10 năm học 2025-2026 TP.HCM – Đề tham khảo Huyện Nhà Bè (Đề 1) – Bài 6
Một người đi xe đạp từ địa điểm $A$ đến địa điểm $B$ với vận tốc $12\text{ km/h}$, sau đó đi tiếp từ $B$ đến địa điểm $C$ với vận tốc $6\text{ km/h}$. Tổng thời gian người đó đi từ $A$ đến $C$ là $75$ phút. Khi quay trở về, người đó đi từ $C$ đến $B$ với vận tốc $8\text{ km/h}$ và từ $B$ về lại $A$ với vận tốc $4\text{ km/h}$, tổng thời gian đi về là $1$ giờ $30$ phút. Tính chiều dài của các quãng đường $AB$ và $BC$.
\loigiai{
Đổi đơn vị thời gian:
$$ 75\text{ phút} = \dfrac{75}{60} = \dfrac{5}{4}\text{ giờ} = 1{,}25\text{ giờ}; \quad 1\text{ giờ } 30\text{ phút} = 1{,}5\text{ giờ} = \dfrac{3}{2}\text{ giờ}. $$
Gọi $x$ và $y$ lần lượt là độ dài của quãng đường $AB$ và quãng đường $BC$ ($x, y > 0$, đơn vị: km).\\
Thời gian người đó đi chiều đi từ $A$ đến $C$ là:
$$ \dfrac{x}{12} + \dfrac{y}{6} = \dfrac{5}{4} \iff x + 2y = 15. $$
Thời gian người đó đi chiều về từ $C$ về $A$ là:
$$ \dfrac{y}{8} + \dfrac{x}{4} = \dfrac{3}{2} \iff 2x + y = 12. $$
Ta có hệ phương trình bậc nhất hai ẩn:
$$ \begin{cases} x + 2y = 15 \\ 2x + y = 12 \end{cases} \iff \begin{cases} 2x + 4y = 30 \\ 2x + y = 12 \end{cases} \iff \begin{cases} 3y = 18 \\ x = 12 - 2y \end{cases} \iff \begin{cases} y = 6\text{ (thoả mãn)} \\ x = 12 - 2 \cdot 6 = 3\text{ (thoả mãn)}. \end{cases} $$
Vậy quãng đường $AB$ dài $3\text{ km}$ và quãng đường $BC$ dài $6\text{ km}$.
}
\end{bt}

\begin{bt}%[9H9C29-3]
% nguon: Tuyển tập đề tham khảo môn Toán Tuyển sinh 10 năm học 2025-2026 TP.HCM – Đề tham khảo Huyện Nhà Bè (Đề 1) – Bài 7
Cho nửa đường tròn $(O)$ có đường kính $AB$. Lấy điểm $M$ tuỳ ý nằm trên nửa đường tròn ($M \ne A, M \ne B$). Tiếp tuyến tại $M$ cắt hai tiếp tuyến tại $A$ và $B$ của nửa đường tròn $(O)$ lần lượt tại các điểm $C$ và $D$.
\begin{enumerate}[a)]
	\item Chứng minh tứ giác $ACMO$ nội tiếp một đường tròn và $\widehat{CAM} = \widehat{ODM}$.
	\item Gọi $P$ là giao điểm của đường thẳng $CD$ và đường thẳng $AB$. Chứng minh $PA \cdot PO = PC \cdot PM$.
	\item Gọi $E$ là giao điểm của $AM$ và $BD$; $F$ là giao điểm của $AC$ và $BM$. Chứng minh ba điểm $E, F, P$ thẳng hàng.
\end{enumerate}
\loigiai{
\begin{enumerate}[a)]
	\item Vì $CA, CM$ là hai tiếp tuyến của nửa đường tròn $(O)$ nên $\widehat{CAO} = 90^\circ$ và $\widehat{CMO} = 90^\circ$.\\
	Tứ giác $ACMO$ có $\widehat{CAO} + \widehat{CMO} = 90^\circ + 90^\circ = 180^\circ$, suy ra tứ giác $ACMO$ nội tiếp đường tròn đường kính $CO$.\\
	Do tứ giác $ACMO$ nội tiếp nên $\widehat{CAM} = \widehat{COM}$ (hai góc nội tiếp cùng chắn cung $CM$).\\
	Mặt khác, theo tính chất của hai tiếp tuyến cắt nhau, tia $OD$ là phân giác của góc $\widehat{MOB}$, suy ra:
	$$ \widehat{ODM} = 90^\circ - \widehat{MOD} = 90^\circ - \dfrac{1}{2}\widehat{MOB} = \dfrac{1}{2}\widehat{MOA} = \widehat{COM}. $$
	Do đó $\widehat{CAM} = \widehat{ODM}$.
	\item Trong tam giác vuông $PCO$ vuông tại $C$, kẻ đường cao $CM \perp PO$ (vì $CD$ là tiếp tuyến tại $M$ của $(O)$ nên $OM \perp CD$ tại $M$, hay tam giác $PMO$ vuông tại $M$).\\
	Áp dụng hệ thức lượng hoặc xét hai tam giác vuông đồng dạng $\Delta PAM$ và $\Delta POC$ (hoặc phương tích từ điểm $P$ đến đường tròn ngoại tiếp tứ giác $ACMO$):\\
	Tứ giác $ACMO$ nội tiếp đường tròn đường kính $CO$, cát tuyến $PAC$ và cát tuyến $PMO$ xuất phát từ điểm $P$ ngoài đường tròn, suy ra:
	$$ PA \cdot PO = PC \cdot PM. $$
	\item Vì $AC \parallel BD$ (cùng vuông góc với đường kính $AB$), áp dụng định lý Ta-lét:
	\begin{itemize}
		\item Trong tam giác $EBD$ có $AC \parallel BD \implies \dfrac{EA}{EM} = \dfrac{AC}{BD}$.
		\item Trong tam giác $FBM$ có $AC \parallel BD \implies \dfrac{FA}{FC} = \dfrac{FB}{FM}$.
	\end{itemize}
	Đoạn thẳng $EF$ song song với $AC$ và $BD$, và đường thẳng nối $EF$ kéo dài cắt $AB$ tại điểm $P$ theo định lý Menelaus cho tam giác $MAB$ hoặc tính chất tứ giác toàn phần tạo bởi các đường thẳng $AM, BM, AC, BD$.\\
	Do đó ba điểm $E, F, P$ thẳng hàng.
\end{enumerate}
}
\end{bt}

\begin{bt}%[9D6H18-2]
% nguon: Tuyển tập đề tham khảo môn Toán Tuyển sinh 10 năm học 2025-2026 TP.HCM – Đề tham khảo Huyện Nhà Bè (Đề 2) – Bài 1
Cho hàm số $y=\dfrac{1}{4}x^2$ có đồ thị là Parabol $(P)$.
\begin{enumerate}[a)]
	\item Vẽ đồ thị $(P)$ trên hệ trục toạ độ $Oxy$.
	\item Tìm toạ độ những điểm thuộc $(P)$ sao cho tung độ gấp $4$ lần hoành độ.
\end{enumerate}
\loigiai{
\begin{enumerate}[a)]
	\item Bảng giá trị tương ứng của $x$ và $y$:
	\begin{center}
	\begin{tabular}{|c|c|c|c|c|c|}
	\hline
	$x$ & $-4$ & $-2$ & $0$ & $2$ & $4$ \\ \hline
	$y=\dfrac{1}{4}x^2$ & $4$ & $1$ & $0$ & $1$ & $4$ \\ \hline
	\end{tabular}
	\end{center}
	Đồ thị $(P)$ là một đường parabol có đỉnh tại gốc toạ độ $O(0;0)$, nhận trục tung $Oy$ làm trục đối xứng và quay bề lõm lên phía trên.
	\item Gọi $M(x_0; y_0)$ là điểm thuộc Parabol $(P)$.\\
	Theo giả thiết, tung độ gấp $4$ lần hoành độ nên $y_0 = 4x_0$.\\
	Vì $M \in (P)$ nên:
	$$ \dfrac{1}{4}x_0^2 = 4x_0 \iff x_0^2 - 16x_0 = 0 \iff x_0(x_0 - 16) = 0 \iff \left[\begin{aligned} x_0 &= 0 \\ x_0 &= 16. \end{aligned}\right. $$
	Với $x_0 = 0 \implies y_0 = 0$, ta được gốc toạ độ $O(0;0)$.\\
	Với $x_0 = 16 \implies y_0 = 4 \cdot 16 = 64$, ta được điểm $M(16; 64)$.\\
	Vậy có hai điểm thoả mãn là $O(0;0)$ và $M(16; 64)$.
\end{enumerate}
}
\end{bt}

\begin{bt}%[9D6H20-2]
% nguon: Tuyển tập đề tham khảo môn Toán Tuyển sinh 10 năm học 2025-2026 TP.HCM – Đề tham khảo Huyện Nhà Bè (Đề 2) – Bài 2
Cho phương trình bậc hai: $2x^2+15x+8=0$.
\begin{enumerate}[a)]
	\item Chứng minh phương trình luôn có hai nghiệm phân biệt $x_1, x_2$.
	\item Không giải phương trình, hãy tính giá trị của biểu thức:
	$$ A = (x_1-x_2)^2 - \dfrac{x_1+x_2}{2}. $$
\end{enumerate}
\loigiai{
\begin{enumerate}[a)]
	\item Phương trình $2x^2+15x+8=0$ có các hệ số $a=2$, $b=15$, $c=8$.\\
	Biệt thức $\Delta$ của phương trình:
	$$ \Delta = 15^2 - 4 \cdot 2 \cdot 8 = 225 - 64 = 161 > 0. $$
	Vì $\Delta > 0$ nên phương trình luôn có hai nghiệm phân biệt $x_1, x_2$.
	\item Theo hệ thức Vi-ét, ta có:
	$$ \begin{cases} S = x_1+x_2 = -\dfrac{15}{2} \\ P = x_1 x_2 = \dfrac{8}{2} = 4. \end{cases} $$
	Biến đổi biểu thức $A$:
	$$ A = (x_1+x_2)^2 - 4x_1 x_2 - \dfrac{x_1+x_2}{2} = S^2 - 4P - \dfrac{S}{2}. $$
	Thay $S = -\dfrac{15}{2}$ và $P = 4$ vào biểu thức $A$, ta được:
	$$ A = \left(-\dfrac{15}{2}\right)^2 - 4 \cdot 4 - \dfrac{-\dfrac{15}{2}}{2} = \dfrac{225}{4} - 16 + \dfrac{15}{4} = \dfrac{240}{4} - 16 = 60 - 16 = 44. $$
\end{enumerate}
}
\end{bt}

\begin{bt}%[9D8H26-3]
% nguon: Tuyển tập đề tham khảo môn Toán Tuyển sinh 10 năm học 2025-2026 TP.HCM – Đề tham khảo Huyện Nhà Bè (Đề 2) – Bài 3
Biểu đồ cột kép biểu diễn số lượng học sinh tham gia giải thi đấu thể thao của một trường trung học cơ sở gồm các khối: Khối 6 (Nam: $6$, Nữ: $5$), Khối 7 (Nam: $8$, Nữ: $7$), Khối 8 (Nam: $10$, Nữ: $9$), Khối 9 (Nam: $11$, Nữ: $8$). Chọn ngẫu nhiên một học sinh tham gia giải thi đấu thể thao của trường đó. Tính xác suất của mỗi biến cố sau:
\begin{enumerate}[a)]
	\item Biến cố $A$: ``Học sinh được chọn là học sinh nam''.
	\item Biến cố $B$: ``Học sinh được chọn là học sinh nữ và không thuộc khối 9''.
\end{enumerate}
\loigiai{
Tổng số học sinh tham gia giải thi đấu thể thao của toàn trường là:
$$ n(\Omega) = (6 + 5) + (8 + 7) + (10 + 9) + (11 + 8) = 11 + 15 + 19 + 19 = 64\text{ (học sinh)}. $$
\begin{enumerate}[a)]
	\item Tổng số học sinh nam tham gia thi đấu là:
	$$ 6 + 8 + 10 + 11 = 35\text{ (học sinh)}. $$
	Xác suất của biến cố $A$ là:
	$$ P(A) = \dfrac{35}{64}. $$
	\item Số học sinh nữ tham gia thi đấu không thuộc khối 9 (thuộc khối 6, 7 hoặc 8) là:
	$$ 5 + 7 + 9 = 21\text{ (học sinh)}. $$
	Xác suất của biến cố $B$ là:
	$$ P(B) = \dfrac{21}{64}. $$
\end{enumerate}
}
\end{bt}

\begin{bt}%[9D1H3-4]
% nguon: Tuyển tập đề tham khảo môn Toán Tuyển sinh 10 năm học 2025-2026 TP.HCM – Đề tham khảo Huyện Nhà Bè (Đề 2) – Bài 4
Bác Năm có một ao nuôi cá hình chữ nhật nằm gọn bên trong một khu đất hình chữ nhật có chiều dài $100\text{ m}$ và chiều rộng $40\text{ m}$. Bác định làm hai lối đi theo chiều rộng cùng có bề rộng là $x\text{ (m)}$ và hai lối đi theo chiều dài cùng có bề rộng là $y\text{ (m)}$.
\begin{enumerate}[a)]
	\item Viết biểu thức thu gọn theo $x$ và $y$ biểu diễn diện tích của ao nuôi cá.
	\item Tính bề rộng lối đi theo chiều dài $y$ (m) biết bề rộng lối đi theo chiều rộng là $x = 1\text{ m}$ và diện tích của ao nuôi cá đạt $3686\text{ m}^2$.
\end{enumerate}
\loigiai{
\begin{enumerate}[a)]
	\item Sau khi chừa lối đi ở hai bên theo chiều rộng và chiều dài:
	\begin{itemize}
		\item Chiều rộng của ao cá là: $40 - 2x\text{ (m)}$ ($0 < x < 20$).
		\item Chiều dài của ao cá là: $100 - 2y\text{ (m)}$ ($0 < y < 50$).
	\end{itemize}
	Biểu thức diện tích ao nuôi cá là:
	$$ S = (100 - 2y)(40 - 2x) = 4000 - 200x - 80y + 4xy\text{ (m}^2\text{)}. $$
	\item Thay $x = 1\text{ m}$ vào công thức kích thước của ao, ta có:
	\begin{itemize}
		\item Chiều rộng ao cá: $40 - 2 \cdot 1 = 38\text{ (m)}$.
		\item Chiều dài ao cá: $100 - 2y\text{ (m)}$.
	\end{itemize}
	Theo đề bài, diện tích ao nuôi cá là $3686\text{ m}^2$, ta có phương trình:
	$$ 38 \cdot (100 - 2y) = 3686 \iff 100 - 2y = \dfrac{3686}{38} = 97 \iff 2y = 3 \iff y = 1{,}5\text{ (m)}. $$
	Vậy bề rộng của lối đi theo chiều dài là $y = 1{,}5\text{ m}$.
\end{enumerate}
}
\end{bt}

\begin{bt}%[9H0V31-3]
% nguon: Tuyển tập đề tham khảo môn Toán Tuyển sinh 10 năm học 2025-2026 TP.HCM – Đề tham khảo Huyện Nhà Bè (Đề 2) – Bài 5
Để làm một chiếc gàu tát nước có dạng hình nón, bác Năm dùng một tấm tôn hình tam giác cân $OMN$ tại $O$ có cạnh bên $OM = ON = 6\text{ dm}$ và góc ở đỉnh $\widehat{MON} = 120^\circ$. Bác xác định trung điểm $H$ của cạnh đáy $MN$, rồi vẽ cung tròn tâm $O$ bán kính $OH$ cắt các cạnh $OM, ON$ lần lượt tại $A, B$. Sau đó bác cắt bỏ phần tôn bên ngoài cung tròn, rồi cuộn phần tôn hình quạt tròn $OAB$ lại sao cho mép $OA$ trùng khít với mép $OB$ để tạo thành chiếc gàu hình nón.
\begin{enumerate}[a)]
	\item Tính diện tích xung quanh của chiếc gàu hình nón (kết quả làm tròn đến $1$ chữ số thập phân).
	\item Hỏi khi múc đầy thì chiếc gàu chứa được bao nhiêu lít nước? (lấy $\pi \approx 3{,}14$ và làm tròn kết quả đến hàng đơn vị).
\end{enumerate}
\loigiai{
\begin{enumerate}[a)]
	\item Tam giác $OMN$ cân tại $O$ có $OM = 6\text{ dm}$, góc $\widehat{MON} = 120^\circ$. Vì $H$ là trung điểm của $MN$ nên $OH \perp MN$ và $OH$ là tia phân giác của $\widehat{MON} \implies \widehat{MOH} = 60^\circ$.\\
	Trong tam giác vuông $OHM$ vuông tại $H$:
	$$ OH = OM \cdot \cos \widehat{MOH} = 6 \cdot \cos 60^\circ = 6 \cdot \dfrac{1}{2} = 3\text{ dm}. $$
	Chiếc gàu được cuộn từ hình quạt tròn $OAB$ có bán kính $R = OH = 3\text{ dm}$ và góc ở tâm là $120^\circ$.\\
	Diện tích xung quanh của chiếc gàu bằng diện tích hình quạt tròn $OAB$:
	$$ S_{xq} = \dfrac{\pi R^2 \cdot 120^\circ}{360^\circ} = \dfrac{\pi \cdot 3^2}{3} = 3\pi \approx 3 \cdot 3{,}14 = 9{,}42 \approx 9{,}4\text{ dm}^2. $$
	\item Độ dài cung tròn $AB$ của hình quạt tròn là:
	$$ l = \dfrac{2\pi R \cdot 120^\circ}{360^\circ} = \dfrac{2\pi \cdot 3}{3} = 2\pi\text{ (dm)}. $$
	Khi cuộn thành hình nón, chu vi đáy của hình nón bằng độ dài cung $AB$:
	$$ 2\pi r = l = 2\pi \implies r = 1\text{ dm}. $$
	Đường sinh của hình nón là bán kính của hình quạt: $l_{\text{sinh}} = R = 3\text{ dm}$.\\
	Chiều cao của hình nón là:
	$$ h = \sqrt{l_{\text{sinh}}^2 - r^2} = \sqrt{3^2 - 1^2} = \sqrt{8} = 2\sqrt{2}\text{ dm}. $$
	Thể tích của chiếc gàu hình nón là:
	$$ V = \dfrac{1}{3}\pi r^2 h = \dfrac{1}{3} \cdot 3{,}14 \cdot 1^2 \cdot 2\sqrt{2} \approx 2{,}96\text{ dm}^3. $$
	Vì $1\text{ dm}^3 = 1\text{ lít}$ nên làm tròn đến hàng đơn vị, chiếc gàu chứa được khoảng $3$ lít nước.
\end{enumerate}
}
\end{bt}

\begin{bt}%[9D1H3-4]
% nguon: Tuyển tập đề tham khảo môn Toán Tuyển sinh 10 năm học 2025-2026 TP.HCM – Đề tham khảo Huyện Nhà Bè (Đề 2) – Bài 6
Vào buổi lễ tốt nghiệp của lớp 9A, cô giáo chủ nhiệm tặng mỗi bạn học sinh nam $2$ cây viết và $1$ quyển tập, còn mỗi bạn học sinh nữ thì được tặng $3$ cây viết và $2$ quyển tập. Cô đã mua từ cửa hàng tổng cộng $124$ cây viết và $75$ quyển tập để phát đủ cho cả lớp. Biết rằng ngày hôm đó tất cả học sinh của lớp 9A đều có mặt dự lễ và cùng chụp chung một tấm ảnh tập thể. Hỏi lớp 9A có tất cả bao nhiêu học sinh?
\loigiai{
Gọi số học sinh nam và số học sinh nữ của lớp 9A lần lượt là $x$ và $y$ ($x, y \in \mathbb{N}^*$, học sinh).\\
Tổng số cây viết cô giáo tặng cho cả lớp là:
$$ 2x + 3y = 124. $$
Tổng số quyển tập cô giáo tặng cho cả lớp là:
$$ x + 2y = 75. $$
Ta có hệ phương trình bậc nhất hai ẩn:
$$ \begin{cases} 2x + 3y = 124 \\ x + 2y = 75 \end{cases} \iff \begin{cases} 2(75 - 2y) + 3y = 124 \\ x = 75 - 2y \end{cases} \iff \begin{cases} 150 - y = 124 \\ x = 75 - 2y \end{cases} \iff \begin{cases} y = 26\text{ (thoả mãn)} \\ x = 75 - 2 \cdot 26 = 23\text{ (thoả mãn)}. \end{cases} $$
Số học sinh nam là $23$ bạn và số học sinh nữ là $26$ bạn.\\
Vậy tổng số học sinh của lớp 9A có trong tấm ảnh tập thể là:
$$ 23 + 26 = 49\text{ (học sinh)}. $$
}
\end{bt}

\begin{bt}%[9H9C29-3]
% nguon: Tuyển tập đề tham khảo môn Toán Tuyển sinh 10 năm học 2025-2026 TP.HCM – Đề tham khảo Huyện Nhà Bè (Đề 2) – Bài 7
Cho tam giác $ABC$ có ba góc nhọn với $\widehat{ABC} = 60^\circ$, $BC = 2a$ và $AB < AC$. Gọi $(O)$ là đường tròn đường kính $BC$ ($O$ là trung điểm cạnh $BC$). Đường tròn $(O)$ cắt các cạnh $AB, AC$ lần lượt tại $D$ và $E$ ($D \ne B, E \ne C$). Đoạn thẳng $BE$ cắt $CD$ tại trực tâm $H$.
\begin{enumerate}[a)]
	\item Chứng minh tứ giác $ADHE$ nội tiếp một đường tròn và xác định tâm $I$ của đường tròn ngoại tiếp tứ giác đó.
	\item Tiếp tuyến tại $C$ của đường tròn $(O)$ cắt đường thẳng $DI$ tại điểm $M$. Chứng minh $HB \cdot DE = HD \cdot BC$ và $OM = 2OC$.
	\item Gọi $F$ là giao điểm của $AH$ và cạnh $BC$. Kẻ $HK$ vuông góc với đoạn thẳng $DF$ tại điểm $K$. Biết độ dài đoạn thẳng $BF = \dfrac{3a}{4}$, hãy tính độ dài đoạn thẳng $HK$ theo $a$.
\end{enumerate}
\loigiai{
\begin{enumerate}[a)]
	\item Vì $D, E$ thuộc đường tròn $(O)$ đường kính $BC$ nên $\widehat{BDC} = 90^\circ$ và $\widehat{BEC} = 90^\circ$.\\
	Suy ra $CD \perp AB$ và $BE \perp AC$, do đó $\widehat{ADH} = 90^\circ$ và $\widehat{AEH} = 90^\circ$.\\
	Tứ giác $ADHE$ có $\widehat{ADH} + \widehat{AEH} = 90^\circ + 90^\circ = 180^\circ$, suy ra tứ giác $ADHE$ nội tiếp đường tròn đường kính $AH$.\\
	Tâm $I$ của đường tròn này là trung điểm của đoạn thẳng $AH$.
	\item Xét tam giác $HDE$ và tam giác $HBC$ có:
	$$ \widehat{DHE} = \widehat{BHC}\text{ (đối đỉnh)}, \quad \widehat{HDE} = \widehat{HBC}\text{ (hai góc nội tiếp cùng chắn cung } EC \text{ của } (O)). $$
	Suy ra $\Delta HDE \backsim \Delta HBC$ (g-g) $\implies \dfrac{HD}{HB} = \dfrac{DE}{BC} \implies HB \cdot DE = HD \cdot BC$.\\
	Vì $I$ là tâm đường tròn ngoại tiếp tam giác vuông $ADH$ nên $ID = IH \implies \Delta IDH$ cân tại $I \implies \widehat{IDH} = \widehat{IHD}$.\\
	Vì $O$ là trung điểm $BC$ nên $OD = OB \implies \Delta OBD$ cân tại $O$ có $\widehat{B} = 60^\circ \implies \Delta OBD$ đều $\implies \widehat{ODB} = 60^\circ$.\\
	Ta có $\widehat{ODI} = \widehat{ODB} + \widehat{BDC} - \widehat{IDH} = 90^\circ$, suy ra $DI$ là tiếp tuyến của đường tròn $(O)$ tại tiếp điểm $D$.\\
	Vì $MD$ và $MC$ là hai tiếp tuyến của $(O)$ nên $MO$ là tia phân giác của góc $\widehat{DOC}$.\\
	Góc ở tâm $\widehat{DOC} = 180^\circ - \widehat{BOD} = 180^\circ - 60^\circ = 120^\circ \implies \widehat{COM} = \dfrac{120^\circ}{2} = 60^\circ$.\\
	Trong tam giác vuông $OCM$ vuông tại $C$:
	$$ \cos \widehat{COM} = \dfrac{OC}{OM} \implies \cos 60^\circ = \dfrac{OC}{OM} = \dfrac{1}{2} \implies OM = 2OC. $$
	\item Vì $H$ là trực tâm của tam giác $ABC$ nên $AF \perp BC$ tại $F$.\\
	Ta có $BC = 2a$ và $BF = \dfrac{3a}{4} \implies CF = BC - BF = 2a - \dfrac{3a}{4} = \dfrac{5a}{4}$.\\
	Trong tam giác vuông $ABF$ vuông tại $F$ có $\widehat{B} = 60^\circ$:
	$$ AF = BF \cdot \tan 60^\circ = \dfrac{3a}{4} \cdot \sqrt{3} = \dfrac{3\sqrt{3}a}{4}. $$
	Trong tam giác $ABC$, hệ thức đường cao cho ta:
	$$ AF \cdot HF = BF \cdot CF \implies HF = \dfrac{BF \cdot CF}{AF} = \dfrac{\dfrac{3a}{4} \cdot \dfrac{5a}{4}}{\dfrac{3\sqrt{3}a}{4}} = \dfrac{5a}{4\sqrt{3}} = \dfrac{5\sqrt{3}a}{12}. $$
	Trong tam giác vuông $BDF$ vuông tại $F$, ta có:
	$$ DF = \sqrt{BF^2 + BD^2 - 2BF \cdot BD \cdot \cos 60^\circ}. $$
	Vì tam giác $OBD$ đều nên $BD = OB = a$, do đó:
	$$ DF = \sqrt{\left(\dfrac{3a}{4}\right)^2 + a^2 - 2 \cdot \dfrac{3a}{4} \cdot a \cdot \dfrac{1}{2}} = \sqrt{\dfrac{9a^2}{16} + a^2 - \dfrac{3a^2}{4}} = \dfrac{\sqrt{13}a}{4}. $$
	Trong tam giác vuông $DFH$, kẻ đường cao $HK \perp DF$, ta có diện tích tam giác $DFH$ bằng $\dfrac{1}{2}HK \cdot DF = \dfrac{1}{2}HF \cdot d(D, AF)$.\\
	Khoảng cách từ $D$ đến $AF$ là độ dài đoạn vuông góc từ $D$ xuống $AF$, bằng $\dfrac{a}{4}$.\\
	Do đó:
	$$ HK = \dfrac{HF \cdot \dfrac{a}{4}}{DF} = \dfrac{\dfrac{5\sqrt{3}a}{12} \cdot \dfrac{a}{4}}{\dfrac{\sqrt{13}a}{4}} = \dfrac{5\sqrt{3}a}{12\sqrt{13}} = \dfrac{5\sqrt{39}}{156}a. $$
\end{enumerate}
}
\end{bt}
